\section{Higher-order collisions: quadratic collapse and shape obstruction}
\label{sec:multiple-collision-main}

The double-pole experiment retains the magnitude of the rate split because a
centred two-point configuration is determined, up to permutation, by its
quadratic dispersion.  For three or more colliding rates this is no longer
true.  The next result identifies the entire local experiment: at the
fourth-root scale it retains the quadratic dispersion and discards the shape
of the splitting configuration.  This yields an efficient test of an
arbitrary isolated multiple collision and, simultaneously, an obstruction to
recovering the local cluster shape.

Fix \(3\le m\le n\).  Define the compact \(m\)-fold collision stratum
\begin{equation}
 \mathcal C_{n,m}(\varepsilon,M,c_0)=
 \left\{(\underbrace{\lambda,\ldots,\lambda}_{m},
 \lambda_{m+1},\ldots,\lambda_n):
 \begin{array}{l}
 \lambda,\lambda_j\in[\varepsilon,M],\quad
 \lambda_{m+1}<\cdots<\lambda_n,\\
 |\lambda_j-\lambda|\ge c_0,\quad
 |\lambda_i-\lambda_j|\ge c_0
 \end{array}\right\},
\label{eq:main-multiple-collision-stratum}
\end{equation}
where the last separation condition ranges over \(m+1\le i<j\le n\).
All conditions involving separated rates are vacuous when \(m=n\).
At a base point write
\(\eta_0=(\lambda,\lambda_{m+1},\ldots,\lambda_n)\).  For bounded
\(u=(u_c,u_{m+1},\ldots,u_n)\) and
\begin{equation}
 \mathcal H_m(L)=\left\{h\in\RR^m:
 \sum_{i=1}^m h_i=0,\ \|h\|\le L\right\},
 \qquad \rho(h)=\frac12\sum_{i=1}^m h_i^2,
\label{eq:main-cluster-dispersion}
\end{equation}
consider the rate multiset
\begin{equation}
 \mathcal R_N(u,h)=
 \left\{\!\left\{
 \lambda+\frac{u_c}{\sqrt N}+\frac{h_i}{N^{1/4}}:1\le i\le m;
 \quad \lambda_j+\frac{u_j}{\sqrt N}:m+1\le j\le n
 \right\}\!\right\}.
\label{eq:main-multiple-local-rates}
\end{equation}
Let \(P_{N,u,h}\) denote the stationary renewal-window law generated by
these rates, and let \(P_{N,0,0}\) denote the law at the collision base.

To define the information, put \(d_i=\theta_i-c\), where
\(c=m^{-1}\sum_{i=1}^m\theta_i\), and use the symmetric cluster chart
\begin{equation}
 \prod_{i=1}^m(x+d_i)
 =x^m-\rho x^{m-2}+\sum_{j=3}^m a_jx^{m-j}.
\label{eq:main-symmetric-cluster-chart}
\end{equation}
The sampled gap mass has an analytic continuation in
\((\eta,\rho,a_3,\ldots,a_m)\) at the collision.  Let
\(s_{\eta,\rho}\) be its score in the coordinates \((\eta,\rho)\), with
the remaining \(a_j\)'s fixed at zero, and set
\begin{equation}
 I_{n,m}(\eta_0)=\frac1{\mu_{\eta_0}}
 \sum_{k\ge0}g_{\eta_0}(k)
 s_{\eta,\rho}(k)s_{\eta,\rho}(k)^\top.
\label{eq:main-multiple-information}
\end{equation}

\begin{theorem}[Quadratic collapse at an isolated multiple pole]
\label{thm:multiple-collision-collapse-main}
Fix \(n,m,\varepsilon,M,c_0\), known \(\Delta\tau>0\), and the serial
positive killed-reset realization, and suppose
\(\mathcal C_{n,m}(\varepsilon,M,c_0)\ne\varnothing\).

\smallskip\noindent\textup{(A) Collapsed local experiment.}
The matrix \(I_{n,m}(\eta_0)\) is uniformly positive definite on the
stratum.  Uniformly over its base points, bounded \(u\), and
\(h\in\mathcal H_m(L)\),
\begin{equation}
 \log\frac{dP_{N,u,h}}{dP_{N,0,0}}
 =\begin{pmatrix}u\\ \rho(h)\end{pmatrix}^{\!\top}\Delta_{N,\eta_0}
 -\frac12\begin{pmatrix}u\\ \rho(h)\end{pmatrix}^{\!\top}
 I_{n,m}(\eta_0)
 \begin{pmatrix}u\\ \rho(h)\end{pmatrix}
 +o_{P_{N,0,0}}(1),
\label{eq:main-multiple-lan}
\end{equation}
where
\(\Delta_{N,\eta_0}\Rightarrow
\mathcal N\{0,I_{n,m}(\eta_0)\}\) uniformly in bounded-Lipschitz
distance.  Hence the local experiment indexed by \(h\) factors through the
single coordinate \(h\mapsto\rho(h)\).

\smallskip\noindent\textup{(B) Efficient test of a multiple collision.}
Partition \(I_{n,m}\) according to \((\eta,\rho)\) and let
\begin{equation}
 J_{n,m}=I_{\rho\rho}
 -I_{\rho\eta}I_{\eta\eta}^{-1}I_{\eta\rho}>0.
\label{eq:main-multiple-schur-complement}
\end{equation}
There is a total measurable residualized-score test of
\(h=0\) against \(h\ne0\) with uniform local asymptotic level
\(\alpha\) and limiting power
\begin{equation}
 1-\Phi_{\rm N}\{z_{1-\alpha}-\rho(h)\sqrt{J_{n,m}(\eta_0)}\}.
\label{eq:main-multiple-power-envelope}
\end{equation}
This is the Gaussian half-line power envelope among tests with uniform local
asymptotic level \(\alpha\).

\smallskip\noindent\textup{(C) Shape obstruction and recovery rate.}
There are \(h,h'\in\mathcal H_m(L)\), for some finite \(L\), such that
\(\rho(h)=\rho(h')\) but their coordinate multisets are distinct.  Holding
\(u=0\),
\begin{equation}
 \|P_{N,0,h}-P_{N,0,h'}\|_{\rm TV}\longrightarrow0.
\label{eq:main-shape-tv-collapse}
\end{equation}
If
\(b=d_{\rm m}(\{\!\{h_1,\ldots,h_m\}\!\},
\{\!\{h'_1,\ldots,h'_m\}\!\})>0\), then, for every \(0<a<b/2\),
\begin{equation}
 \liminf_{N\to\infty}\inf_{\widehat{\mathcal R}_N}
 \max_{\xi\in\{h,h'\}}P_{N,0,\xi}\left\{
 N^{1/4}d_{\rm m}(\widehat{\mathcal R}_N,
 \mathcal R_N(0,\xi))>a\right\}\ge\frac12.
\label{eq:main-shape-minimax-lower}
\end{equation}
Conversely, a measurable estimator is uniformly
\(O_P(N^{-1/4})\) over bounded \(u\) and \(h\).  Thus \(N^{-1/4}\) is the
optimal local multiset-recovery rate, but for \(m\ge3\) the scaled cluster
shape is not consistently estimable.
\end{theorem}

\begin{proof}
Let \(B_\eta(s)=\prod_{j=m+1}^n\theta_j/(s+\theta_j)\), with the empty
product equal to one.  In the chart
\eqref{eq:main-symmetric-cluster-chart}, the cluster factor in the Laplace
transform of the holding-time sum is
\begin{equation}
 \frac{c^m-\rho c^{m-2}+\sum_{j=3}^m a_jc^{m-j}}
 {(s+c)^m-\rho(s+c)^{m-2}
   +\sum_{j=3}^m a_j(s+c)^{m-j}}.
\label{eq:main-multiple-cluster-transform}
\end{equation}
This rational expression supplies the asserted analytic continuation.  If
\(H\) denotes the Laplace transform of the continuous survival function,
differentiation at the collision gives
\begin{equation}
 \partial_\rho H(s)
 =c^{m-2}B_\eta(s)\frac{s+2c}{(s+c)^{m+2}},
 \qquad
 \partial_cH(s)=-mc^{m-1}B_\eta(s)\frac1{(s+c)^{m+1}}.
\label{eq:main-multiple-transform-derivatives}
\end{equation}
The first derivative has an exact pole of order \(m+2\) at \(-c\), while
the centre derivative has order \(m+1\).  A separated-rate derivative has
an exact order-two pole at its own pole.  Inverse Laplace transformation and
sampling therefore give unique leading components
\begin{equation}
 k^{m+1}e^{-c\Delta\tau k},\qquad
 k^me^{-c\Delta\tau k},\qquad
 ke^{-\theta_j\Delta\tau k},
\label{eq:main-multiple-distinguishing-components}
\end{equation}
respectively.  The exponential-polynomial independence argument in
\Cref{prop:sampled-double-pole-nondegeneracy-main} proves linear independence
of the \((\eta,\rho)\) scores.  Uniform separation, dominated summation, and
compactness then prove uniform positive definiteness of
\eqref{eq:main-multiple-information}.

We next establish the local expansion, including the absence of higher
cluster-shape coordinates.  If \(e_j(h)\) is the \(j\)th elementary
symmetric polynomial, then the local split in
\eqref{eq:main-multiple-local-rates} has
\begin{equation}
 \rho_N=N^{-1/2}\rho(h),\qquad
 a_{j,N}=N^{-j/4}e_j(h),\quad 3\le j\le m.
\label{eq:main-multiple-coordinate-scales}
\end{equation}
In particular, every \(a_{j,N}=o(N^{-1/2})\).

For completeness, the required quadratic-mean remainder follows directly
from the complete holding times.  At a collision base, their likelihood
ratio for a centred split \(d=(d_1,\ldots,d_m)\) contains the factor
\begin{equation}
 L_d(x)=\prod_{i=1}^m(1+d_i/c)
 \exp\left(-\sum_{i=1}^m d_ix_i\right).
\label{eq:main-complete-split-likelihood}
\end{equation}
Conditional on \(Y=\sum_{i=1}^mX_i\), the base law of
\((X_1,\ldots,X_m)\) is exchangeable.  Hence the conditional expectation
of the linear term in \(d\) vanishes when \(\sum_i d_i=0\).  Its quadratic
term is an exchangeable quadratic form on that hyperplane and is therefore
a scalar multiple of \(\sum_i d_i^2=2\rho\); comparison with
\eqref{eq:main-multiple-transform-derivatives} identifies it as the
\(\rho\)-score.  Taylor's formula in
\eqref{eq:main-complete-split-likelihood} bounds the remaining conditional
term by
\(C\|d\|^3(1+Y)^3e^{C\|d\|Y}\).  The common exponential moment makes this
bound square integrable uniformly after conditioning on a sampling bin or a
tail event.  With \(d=N^{-1/4}h\), it is \(O(N^{-3/4})\) in relative
\(L^2\), uniformly for bounded \(h\).  The ordinary root-\(N\) expansion in
\(\eta\), including its cross-remainder with \(d\), therefore gives
\begin{equation}
 \sqrt{g_{N,u,h}/g_0}
 =1+\frac1{2\sqrt N}
 \begin{pmatrix}u\\\rho(h)\end{pmatrix}^{\!\top}
 s_{\eta,\rho}+o_{L^2(g_0)}(N^{-1/2}),
\label{eq:main-multiple-dqm}
\end{equation}
uniformly on the stated sets.  The identical conditional argument for a
tail event gives the equilibrium residual-life expansion and a common
exponential envelope.

Consequently the hypotheses of
\Cref{thm:stationary-renewal-palm-equivalence-main} hold uniformly for these
arrays.  Its Palm sample has size asymptotic to \(N/\mu_{\eta_0}\).
The i.i.d. DQM theorem applied to
\eqref{eq:main-multiple-dqm}, followed by the two-sided deficiency transfer,
proves \eqref{eq:main-multiple-lan}.  It also proves the factorization through
\(\rho(h)\): the coordinates \(a_3,\ldots,a_m\) lie below local statistical
resolution.

For the test, residualize the \(\rho\)-score against the \(\eta\)-scores.
The pole argument makes its variance \(J_{n,m}\) strictly positive.  The
leading Hankel recurrence and an isolating contour around the \(m\)-root
cluster give a measurable root-\(N\) null fit of its centre and of the
separated rates, exactly as in the finite contour atlas used in
\Cref{est:serial-collision-recurrence-main}.  Orthogonality of the efficient
score to the nuisance scores, polynomial score envelopes, and the stopped
renewal LLN show that inserting this fit and empirical information changes
the standardized efficient score by \(o_P(1)\).  Le Cam's third lemma gives
mean \(\rho(h)\sqrt{J_{n,m}}\) and unit variance.  The
Neyman--Pearson lemma in the residual Gaussian shift proves
\eqref{eq:main-multiple-power-envelope} and its envelope assertion.

It remains to prove the stronger shape statement.  For \(m=3\), take, for
any sufficiently small \(r>0\),
\begin{equation}
 h=(r,-r,0),\qquad
 h'=(r/\sqrt3,r/\sqrt3,-2r/\sqrt3),
\label{eq:main-equal-dispersion-shapes}
\end{equation}
and append zeros for \(m>3\).  Both vectors are centred and have
\(\rho=r^2\), but their multisets differ.  Thus \(b>0\).
Because their \(\rho\)-coordinates agree, the relative \(L^2\) expansion
above starts at order \(N^{-3/4}\), so
\begin{equation}
 H^2(g_{N,0,h},g_{N,0,h'})=O(N^{-3/2}),\qquad
 H^2(e_{N,0,h},e_{N,0,h'})=O(N^{-3/2}).
\label{eq:main-equal-dispersion-hellinger}
\end{equation}
For the undershot Palm sample, Hellinger tensorization makes the squared
distance \(O(N^{-1/2})\).  Apply the common reverse kernel from
\Cref{thm:stationary-renewal-palm-equivalence-main} to the two Palm laws.
Its two approximation errors vanish uniformly, while total variation
contracts under a Markov kernel.  This proves
\eqref{eq:main-shape-tv-collapse}.

The two balls of radius \(aN^{-1/4}\) around the rate multisets generated by
\(h\) and \(h'\) are disjoint for all sufficiently large \(N\).  Assigning
an estimator to the nearer candidate converts simultaneous estimation
success into a test.  The equal-prior testing error is at least
\((1-\|P_{N,0,h}-P_{N,0,h'}\|_{\rm TV})/2\), which proves
\eqref{eq:main-shape-minimax-lower}.

Finally, estimate the cluster centre and the separated roots by the same
isolating-contour recurrence fit and return \(m\) copies of the fitted
centre.  The fit is root-\(N\), whereas every true cluster rate differs from
its centre by at most \(LN^{-1/4}\).  Uniform separation fixes the matching,
and hence this total measurable estimator is uniformly
\(O_P(N^{-1/4})\).  Together with
\eqref{eq:main-shape-minimax-lower}, this proves the last assertion.
\end{proof}
